% Shared exercise chunk: l4-sharing
\begin{exercisechunk}
\item \textbf{Estimating shared parameters.\suppmark}
\label[exercise]{ex:sharing}
\exerciseprofile{3}{3}{3}{\exproof\quad\excalculation\quad\exinterpretation}{%
Show why an unsmoothed empirical Bernoulli estimate has infinite expected KL
and prove a finite-sample bound for add-one smoothing.}
Use the independent-bit models and pooled estimator from
\cref{sec:parameter-sharing}. All expectations below average over the
training observations. Fix \(0<a<1/2\).
Let \(B\sim\operatorname{Binomial}(L,p)\), with \(L\ge1\).
For \(0<p<1\), show that the ordinary empirical frequency
\(\bar p=B/L\) satisfies \(\E[\kl(p,\bar p)]=\infty\).
Explain why adding one count of each bit gives
\(\widehat p=(B+1)/(L+2)\), and why this assigns positive probability
to both bits. To obtain a finite-sample bound for this choice, prove
that, for \(0<p<1\),
\[
 \E\!\left[\frac1{B+1}\right]
 =\frac{1-(1-p)^{L+1}}{(L+1)p}.
\]
Apply the same identity to \(L-B\), then use Jensen's inequality
for the logarithm to show
\[
 \E[\kl(p,\widehat p)]
 \le\log\frac{L+2}{L+1}
 \le\frac1{L+1}.
\]
Verify the bound directly for \(p=0,1\).
\begin{hints}
Use \(1/(B+1)=\int_0^1u^B\,du\).
\end{hints}
\end{exercisechunk}
